\documentclass[10pt,a4paper]{amsart}
\usepackage[margin=19mm]{geometry}
\usepackage{amsmath,amssymb,mathtools}
\usepackage[scr=rsfs,cal=euler]{mathalfa}
\usepackage{xfrac}
\usepackage[unicode,hidelinks,bookmarks=false]{hyperref}
\newcommand{\ie}{i.e.\ }
\newcommand{\cf}{cf.\ }
\newcommand{\resp}{respectively\ }
\newcommand{\Subsection}{\subsection}
\providecommand{\nobreakdash}{\nobreak\hbox{-}}
\setlength{\emergencystretch}{2em}
\multlinegap0pt
\usepackage{bm}
\usepackage{bbm}
\usepackage{dsfont}
\usepackage{xspace}
\usepackage{cancel}
\usepackage{mathtools}
\usepackage{amsthm}
\makeatletter
\@ifundefined{thm}{\newtheorem{thm}{Theorem}}{}
\makeatother
\newcommand\BA{\mathbf A}
\newcommand\Bc{\mathbf c}
\newcommand\Bd{\mathbf d}
\newcommand\Be{\mathbf e}
\newcommand\Bh{\mathbf h}
\newcommand\ul{\underline}
\newcommand\weit{\operatorname{wt}}
\newcommand{\NN}{\mathbb N}
\renewcommand\b{\beta}
\title[Algorithm for computing canonical bases and foldings of quan]{Algorithm for computing canonical bases and foldings of quantum groups}
\author{Toshiaki Shoji, Zhiping Zhou}
\date{Source: arXiv:2506.00793v1 (CC BY 4.0)}
\begin{document}
\maketitle

\noindent\emph{Source and licence.} Algorithm for computing canonical bases and foldings of quantum groups. Version arXiv:2506.00793v1 (CC BY 4.0), distributed under the Creative Commons Attribution 4.0 International licence, as recorded in the arXiv licence metadata for this exact version and shown on the abstract page https://arxiv.org/abs/2506.00793v1. Full rights notice: licenses/arxiv-2506.00793-CC-BY-4.0.txt.\par\smallskip

\noindent\textit{Editorial scope.} Counted unit: the Thm whose source label is \texttt{None} in the article (source0.tex lines 2093--2101, source label \texttt{None}), delivered together with the article's own complete original proof of it (source0.tex lines 2103--2202, one \texttt{proof} environment of 100 lines). No mathematics is written, repaired or re-derived: statement and proof are the article's own text line for line. Required context retained uncounted: none. Every definition, display and label of those lines is retained unchanged. Omitted from this package and not redistributed: 1--2092, 2102--2102, 2203--3980. Nothing inside a retained excerpt is omitted or altered: every retained body is delivered exactly as the article's lines are written, so no omission receipt of any kind accompanies this package. Citations: the retained excerpts carry no citation command at all, so no bibliography entry of the article is transcribed and no reference is redistributed. Reference adaptation: none was needed and none was made; every reference inside the delivered unit resolves inside it, so no unresolved reference or citation is produced. Printed slips kept literal: no printed word, symbol, hypothesis or number of the counted statement or of its proof was altered. Layout and class: the article's own class is \texttt{amsart} with packages including latexsym, amssymb, amscd, mathrsfs, bbm, CJK, xy, tikz-cd; the wrapper is the lane's pinned \texttt{amsart} editorial wrapper at 10pt on A4, which restates the theorem-like environments the retained text needs and adds the article's own macro definitions that the retained text needs (\texttt{\textbackslash BA}, \texttt{\textbackslash Bc}, \texttt{\textbackslash Bd}, \texttt{\textbackslash Be}, \texttt{\textbackslash Bh}, \texttt{\textbackslash NN}, \texttt{\textbackslash b}, \texttt{\textbackslash ul}, \texttt{\textbackslash weit}), with extra wrapper packages (bm, bbm, dsfont, xspace, cancel, mathtools). The delivered numbering is the unit's own. Assets: no retained span contains a figure environment, a graphics-inclusion command, a PDF-page inclusion, a table or a TikZ picture; no image file is copied into this package, the package declares no asset dependency and no image file is redistributed with it. Licence: version arXiv:2506.00793v1 of this article is distributed under the Creative Commons Attribution 4.0 International licence, as recorded in the arXiv licence metadata for this exact version and shown on the abstract page https://arxiv.org/abs/2506.00793v1. A full rights notice is delivered at licenses/arxiv-2506.00793-CC-BY-4.0.txt. Characters: every retained excerpt is pure ASCII, so the delivered engine, XeLaTeX with its default Latin Modern text font, reports no missing character.
\begin{thm}  %%%%   Theorem 4.16
For each $\Bc \in \NN^{\ul N}$, $F((\ul\Bc))$ can be written as
\begin{equation*}  
F((\ul\Bc)) = L(\ul\Bc, \ul\Bh)
     + \sum_{\ul\Bc' > \ul\Bc}h_{\ul\Bc',\ul\Bc}L(\ul\Bc', \ul\Bh)
  \quad\text{ with }\quad h_{\ul\Bc',\ul\Bc} \in \BA,  
\end{equation*}
where $\Bc'$ satisfies the relation $\weit(L(\Bc',\ul\Bh)) = \weit(L(\ul\Bc, \ul\Bh))$.
\end{thm}  

\begin{proof}
By induction on $k$, we shall prove
\begin{equation*}
\tag{4.16.1}  
  \ul F(\ul \Bd^1)\ul F(\ul \Bd^2)\cdots \ul F(\ul \Bd^k)
  = L(\ul \Bc_{\le k}, \ul\Bh) + \sum_{\ul\Bc' > \ul\Bc_{\le k}}
          h_{\ul\Bc', \ul\Bc_{\le k}}L(\ul\Bc',\ul\Bh),
\end{equation*}
where $\ul\Bc_{\le k} = (\ul c_1, \dots, \ul c_k, 0, \dots, 0)$, and
$h_{\ul\Bc', \ul\Bc_{\le k}} \in \BA$.  
The formula (4.16.1) certainly holds for $k = 1$ by Lemma 4.15. If it holds for $k = \ul N$,
it gives the assertion of the theorem. 
We assume that (4.16.1) hods for $k$.  Then by Lemma 4.15, we have

\begin{align*}
\tag{4.16.2}
\ul F&(\ul \Bd^1)\ul F(\ul \Bd^2)\cdots \ul F(\ul \Bd^k)\ul F(\ul \Bd^{k+1})  \\
&= \Bigl(L(\ul\Bc_{\le k}, \ul\Bh) + \sum_{\ul\Bc' > \ul\Bc_{\le k}}
        h_{\ul\Bc',\ul\Bc_{\le k}}L(\ul\Bc', \ul\Bh)\Bigr)
        \Bigl(L(\ul\Bc_{k+1}, \ul\Bh) + \sum_{\ul\Bc'' > \ul\Bc_{k+1}}
             a_{\ul\Bc'',\ul\Bc_{k+1}}L(\ul\Bc'',\ul\Bh)\Bigr).
\end{align*}

\par
We consider the relation $\ul\Bc' > \ul\Bc_{\le k}$ for a fixed $k$.
This condition is given by either $(\ul c'_1, \dots, \ul c'_k) > (\ul c_1, \dots, \ul c_k)$ 
or $(\ul c'_1, \dots, \ul c'_k) = (\ul c_1, \dots, \ul c_k)$ with $\ul c'_{\ell} > 0$
for some $\ell > k$.
But since $L(\ul\Bc',\ul\Bh)$ has the same weight as $L(\ul\Bc_{\le k}, \ul\Bh)$,
which is equal to $\ul c_1\ul\b_1 + \cdots + \ul c_k\ul\b_k$, the latter case does not occur,
and we must have $(\ul c_1', \dots, \ul c'_k) > (\ul c_1, \dots, \ul c_k)$. 
\par
We also consider the relation $\ul\Bc'' > \ul\Bc_{k+1}$ for a fixed $k$.
This condition is given by
either $(\ul c''_1, \dots, \ul c''_{k+1}) > (0, \dots, 0, \ul c_{k+1})$ or
$(\ul c''_1, \dots, \ul c''_{k+1}) =  (0, \dots, 0, \ul c_{k+1})$ with $\ul c''_{\ell} > 0$
for some $\ell > {k+1}$.
But since $L(\ul\Bc'', \ul\Bh)$ has the same weight as
$\ul F_{\ul\b_{k+1}}^{(\ul c_{k+1})}$,
which is equal to $\ul c_{k+1}\ul\b_{k+1}$.
Hence the latter case does not occur, and in the former
case, either $\ul c''_{\ell} > 0$ for some $\ell < k+1$, or
$\ul c''_1 = \cdots = \ul c''_k = 0$ and $\ul c''_{k+1} > \ul c_{k+1}$.  
\par
In the following, we show each factor $L(\ul\Bd, \ul\Bh)L(\ul\Bd',\ul\Bh)$ in (4.16.2)
is a linear combination of $L(\ul\Be, \ul\Bh)$ such that $\ul\Be > \ul \Bc_{\le k+1}$.
\par\medskip
(a) First note that
$L(\ul\Bc_{\le k}, \ul\Bh) L(\ul\Bc_{k+1}, \ul\Bh) = L(\ul\Bc_{\le k+1}, \ul\Bh)$.
This is clear from the definition. 
\par\medskip
(b) Next consider $L(\ul\Bc',\ul\Bh)L(\ul c_{k+1}, \ul\Bh)$ such that
$\ul\Bc' > \ul\Bc_{\le k}$. 
By the above remark, we have $(\ul c'_1, \dots, \ul c'_k) > (\ul c_1, \dots, \ul c_k)$.
Moreover, 
$(\ul F_{\ul\b_{k+1}}^{(\ul c'_{k+1})}\cdots \ul F_{\ul\b_{\ul N}}^{(\ul c'_{\ul N})})
\ul F_{\ul\b_{k+1}}^{(\ul c_{k+1})}$
is a linear combination of $L(\ul\Bd, \ul\Bh)$ such that $s(\ul\Bd) \subset [k+1, \ul N]$
(see 2.4 for the definition $s(\ul\Bd)$).     
It follows that $L(\ul\Bc',\ul\Bh)\ul F_{\ul\b_{k+1}}^{(\ul c_{k+1})}$
is a linear combination of $L(\ul\Be, \ul\Bh)$, where $\ul\Be = (\ul e_1, \dots, \ul e_{\ul N})$
is such that $\ul e_{\ell} = \ul c'_{\ell}$ for $\ell = 1, \dots, k$.
In particular, $\ul\Be > \ul\Bc_{\le k+1}$.  Thus $L(\ul\Bc',\ul\Bh)L(\ul\Bc_{k+1}, \ul\Bh)$
is a linear combination of $L(\ul\Be, \ul\Bh)$ such that $\ul\Be > \ul\Bc_{\le k+1}$.

\par\medskip
(c) \ Next consider $L(\ul\Bc', \ul\Bh)L(\ul \Bc'',\ul\Bh)$, where
$\ul\Bc' > \ul\Bc_{\le k}$ and $\ul\Bc'' > \ul\Bc_{k+1}$.
By the above discussion, we have $(\ul c_1', \dots, \ul c_k') > (c_1, \dots, c_k)$.
In particular, we have $\ul\Bc' > \ul\Bc_{\le k+1}$. 
By Proposition 2.5, $L(\ul\Bc', \ul\Bh)L(\ul \Bc'',\ul\Bh)$ is a linear combination of
$L(\ul\Be, \ul\Bh)$ such that $\ul\Be \ge \ul\Bc'$.  Hence $\ul\Be > \ul\Bc_{\le k+1}$.  
Thus $L(\ul\Bc', \ul\Bh)L(\ul \Bc'',\ul\Bh)$ is a linear combination of $L(\ul\Be,\ul\Bh)$
such that $\ul\Be > \ul\Bc_{\le k+1}$. 

\par\medskip
(d) \ Finally consider $L(\ul\Bc_{\le k}, \ul\Bh)L(\ul\Bc'',\ul\Bh)$, where
$\ul\Bc''> \ul\Bc_{k+1}$. Then by the above remark,
either $\ul c''_{\ell} > 0$ for some $\ell\le k$, or $\ul c_1'' = \cdots = \ul c''_k = 0$
and $\ul c''_{k+1} > \ul c_{k+1}$.
Set $\ul\Bc''_{\le k} = (\ul c''_1, \dots, \ul c''_k, 0, \dots, 0)$.  
By Proposition 2.5, $L(\ul\Bc_{\le k}, \ul\Bh)L(\ul \Bc''_{\le k}, \ul\Bh)$ is
a linear combination of $L(\ul\Be, \ul\Bh)$ such that $\ul\Be \ge \ul\Bc_{\le k}$.
Moreover, the support of $\ul\Be$ is contained in $[1,k]$.
If there exists $1 \le \ell \le k$ such that $\ul e_{\ell} > \ul c_{\ell}$, then
we have $\ul\Be > \ul\Bc_{\le k+1}$.  Hence we may assume that
$(\ul e_1, \dots, \ul e_k) = (\ul c_1, \dots, \ul c_k)$. 
But since the weight of $L(\ul\Be, \ul\Bh)$ coincides with
$\weit(L(\ul\Bc_{\le k}, \ul\Bh)) + \weit(L(\ul\Bc''_{\le k}, \ul\Bh))$, 
we must have $\weit(L(\ul\Bc''_{\le k}, \ul\Bh)) = 0$, and $\ul c_1'' = \cdots = \ul c_k'' = 0$.
In particular, $\ul c''_{k+1} > \ul c_{k+1}$ by the above remark. 
Since $L(\ul\Bc'',\ul\Bh) = F_{\ul\b_{k+1}}^{(\ul c''_{k+1})}\cdots
            F_{\ul\b_{\ul N}}^{(\ul c''_{\ul N})}$,
$L(\ul\Bc_{\le k}, \ul\Bh)L(\ul\Bc'', \ul\Bh)$ coincides with
$L(\ul\Be, \ul\Bh)$ such that $\ul\Be > \ul\Bc_{\le k+1}$. 
Thus $L(\ul\Bc_{\le k}, \ul\Bh)L(\ul\Bc'',\ul\Bh)$ is a linear combination of
$L(\ul\Be, \ul\Bh)$ such that $\ul\Be > \ul \Bc_{\le k+1}$.  
\par\medskip
By (a) $\sim$ (d), (4.16.1) holds for $k +1$.  The theorem is proved. 
\end{proof}  

\end{document}
